\section{Keeping the complete question inside a bounded request}
\label{sec:lossless-context}

The authentication circular exposes a practical obstacle to the whole-source
check. Its retained source contains 263 units. A processing exercise assigned a
synthetic concern to every unit, completed the smaller coverage requests, and
then assembled their concerns for reconciliation. The final request occupied
382,396 bytes against a configured limit of 200,000. The source had not been
lost: the unfinished step was the comparison of all the proposed concerns with
the source as a whole. Finishing each smaller request could not answer that
global question. In particular, an exception recorded near the end could still
change a proposed duty recorded near the beginning.

This failure calls for a distinction between repeated representation and new
information. The coverage records occupied 221,974 bytes. They repeated field
names, source quotations, unit identifiers and some identical explanations.
The source itself remained in the request, so an exact quotation could often be
located there without printing it again. Removing that repetition has a
different effect from replacing several concerns with a summary: an exact
reference can recover every original word, whereas a summary requires a further
interpretation of what matters. The capacity repair therefore uses readable
tables and exact references, and checks that they reconstruct the prior request.

\subsection{What the shorter representation says}

The complete source, the proposed interpretation, every question definition and
the response schema remain literal text in the request. A coverage table
declares its columns once: the unit, its proposed disposition, related questions,
related units, evidence and rationale. Each row supplies one value for each
declared column. A separate concern table retains the concern identifier, related
questions and units, evidence, explanation and original identifier from its
smaller coverage request. Keeping both identifiers preserves the connection
between a global concern and the earlier response that introduced it.

The source units have a fixed order. A unit number in a table refers to that
position in the unchanged source list. An evidence entry contains three
integers: the unit number and the beginning and end positions of the quoted
text. Positions count Unicode codepoints, starting at zero; the ending position
is excluded. For example, in the string ``A\,B\,C'', positions zero through
one select the first character. The implementation also exercises non-ASCII
characters, combining marks, line breaks and non-breaking spaces. It never
normalizes a quotation or uses the length of its UTF-8 encoding as a character
position. A quotation must occur exactly once in its unit, preserving the
pre-existing source-evidence rule.

Rationales and explanations receive a similarly exact treatment. A list stores
each distinct complete string, and the row identifies its position in that list.
Two different explanations therefore remain different even when they concern
the same source passage. A long explanation is retained in full. Questions and
concern identifiers stay literal, so the recipient can follow them without
inferring which question an abbreviated label was meant to denote. The
source-span selection table uses declared columns too; its literal span
identifiers remain available for the response. The response still selects these
identifiers under the existing schema, and the source-reference resolver still
checks the selected quotations against the exact source.

\subsection{The check that permits dispatch}

Let $R$ denote the complete original request after its source-span table has
been prepared. Let $E(R)$ be the readable table representation, and let $D$ be
the deterministic decoder. The system retains $R$ locally before applying the
encoder. Its dispatch condition is equality of the canonical JSON bytes:
\begin{equation}
 \operatorname{bytes}\bigl(D(E(R))\bigr)
   = \operatorname{bytes}(R).
 \label{eq:lossless-request-roundtrip}
\end{equation}
Here canonical JSON fixes the representation of the same JSON value for the
comparison. The check covers the entire request, including the unchanged source,
question definitions, coverage hashes, response schema and instructions. It
therefore compares more than the number of reconstructed rows.

The inverse is straightforward to inspect. A row recovers its named fields by
the declared column order. A unit index recovers the original identifier from
the source list. A quotation recovers the indicated contiguous substring, and
an explanation index recovers its complete string. Applying these operations to
every row restores the two record lists in their original order. The decoder
also rejects unknown columns, wrong row lengths, invalid indices, ambiguous
quotations and unused explanation strings. It rebuilds and verifies the
source-reference table, then checks the reconstructed request against the hash
held with the original request. The hash printed inside the transmitted object
is not accepted as its own independent authenticity evidence.

This is an executed preservation check on each request. It is not a Lean proof
of the Python implementation or a proof that a model understands tables. The
distinction matters because a model could receive every word and still attach
an exception to the wrong duty. The later reconciliation therefore remains
subject to the same question, evidence and uncertainty checks. It must account
for every concern and every question; a revised interpretation invalidates all
its previous coverage and requires that coverage to be performed again.

\subsection{What the executed capacity result establishes}

The retained 382,396-byte request becomes 175,016 bytes under this encoding,
leaving 24,984 bytes below the unchanged limit. Decoding reproduces the original
request exactly. All 263 source units and all 263 concerns survive. An offline
integration run then processes the same full source, recreates the same coverage
records, and completes the final reconciliation using a synthetic response
provider. Its final dispositions retain all concerns as unresolved. This
exercise demonstrates that the whole-context control flow can finish within
the byte bound; it supplies no legal answer to the authentication circular.

The first integration diagnostic also exposed an instructive harness mistake.
Its synthetic respondent described an unencoded reading as an executable
proposal. The validator rejected that response. The repair changed the
synthetic response to retain uncertainty, preserving both the failed diagnostic
and the rejection rule. Changing the rule to make the exercise succeed would
have concealed the same error that the interpretation procedures are intended
to detect.

A second capacity challenge gives each unit a long, unique rationale and each
concern additional evidence. Its exact representation exceeds the same limit.
The system retains that request and its successful round-trip check, then stops
before provider dispatch. This outcome is necessary: an encoding cannot promise
to fit an arbitrarily large collection of different explanations into a fixed
request. A future source that reaches this limit needs a separately justified
capacity or reasoning strategy. Dropping qualifications would change the
question being investigated.

The new encoding is an explicit optional protocol,
\texttt{readable-tables.v1}. Existing recorded routes retain their identities;
a changed encoding cannot reuse a journal as though its request protocol were
unchanged. No live model call was used in this capacity exercise, and the
exhausted original allowance and issue histories remain intact. Model
comprehension of this representation is untested. The subsequent executable
quotation and date-proof extensions are described in
Section~\ref{sec:executable-dates}. Unresolved incorporated sources and an
observed prospective evaluation window remain open.
The flow chart on page~\pageref{map:lossless-context} connects the preservation
check to those continuing uncertainties.
